\SourceLeaf{349}

It follows at once that
\[
(t-a)(d-a) \quad \text{and} \quad (b-a)(c-a)
\]
have the same clise or that the angles
\[
\text{between } AT, AC \quad \text{and} \quad \text{between } AB, AD
\]
are equal.

\Needspace{8\baselineskip}We can also put $t$ in the following elegant form
\[
t = \dfrac{\left(m-\dfrac{n}{\beta\gamma}\right)
  \left(m-\dfrac{n}{\alpha\gamma}\right)
  \left(m-\dfrac{n}{\alpha\beta}\right)}
  {\left(\delta - \dfrac{1}{\alpha\beta\gamma}\right) m n}
  + \left(\left(\dfrac{1}{\alpha}+\dfrac{1}{\beta}+\dfrac{1}{\gamma}\right) m
  - \dfrac{1}{\alpha\beta\gamma} n\right) \dfrac{n}{m}
\]
or in this one
\[
t = \dfrac{(\beta\gamma m - n)(m - \beta\delta n)(m-\gamma\delta n)}
  {\alpha\beta\gamma\delta\left(\beta\gamma\delta - \dfrac{1}{\alpha}\right) m n}
  + \dfrac{\beta\gamma m m - m n + (\beta+\gamma)\delta n n}{\beta\gamma\delta n}.
\]

\Needspace{10\baselineskip}\begin{center}
[2.]
\end{center}

The centre of a circle through three points $A$, $B$, $C$ of the ellipse
is
\[
\dfrac{\left\{\alpha\beta\gamma +
  \left(\dfrac{1}{\alpha}+\dfrac{1}{\beta}+\dfrac{1}{\gamma}\right)\right\} m
  - \left\{(\alpha+\beta+\gamma) + \dfrac{1}{\alpha\beta\gamma}\right\} n}
  {\dfrac{m}{n} - \dfrac{n}{m}},
\]
and its radius equals the modulus of \[
\dfrac{(m-\beta\gamma n)(m-\alpha\gamma n)(m-\alpha\beta n)}{mm - nn}
\]
or (which is the same thing) of
\[
\dfrac{\left(m-\dfrac{n}{\beta\gamma}\right)\left(m-\dfrac{n}{\alpha\gamma}\right)
  \left(m-\dfrac{n}{\alpha\beta}\right)}{mm-nn}.
\]

The centre of the circle just mentioned can also be expressed as follows if, for
brevity, the equations
\[
\beta\gamma = \cos 2\mathfrak{A} + i\sin 2\mathfrak{A}
\]
\[
\alpha\gamma = \cos 2\mathfrak{B} + i\sin 2\mathfrak{B}
\]
\[
\alpha\beta = \cos 2\mathfrak{C} + i\sin 2\mathfrak{C}
\]%
